\section{Three small modes on a four-node toric transition}\label{sec:cy-x9}

We now use the original-side nodal partial resolution $X'_9$ constructed
in CY~VI \cite[Section~11.1 and Appendix~G.1]{wuebben-cy-mirror}.
That paper supplies its four nodes, a projective small resolution and
the integral exceptional-curve relation. The deformation results of
CY~V supply a prescribed first-order smoothing tangent and its
integrability. We specify an exact polarization, verify the remaining
family hypotheses, and compute three spectral coefficients from the
full quotient harmonic norm. This connects the earlier geometric
construction to Theorem~\ref{thm:matrix}.

The notation distinguishes three spaces. The original $X_9$ of CY~V
has two nodes and a degree-seven del Pezzo cone; it does not satisfy
our nodal hypotheses. Its partial resolution $X'_9$ has four nodes.
The four-node mirror studied elsewhere in CY~VI is a different
threefold. Only $X'_9$ is used in this section.

\subsection{The earlier nodal model and its polarization}
The eleven rays $v_0,\ldots,v_{10}$ of the original-side ambient fan
are the columns of
\begin{equation}\label{eq:x9-rays}
 \begin{pmatrix}
 1&0&-1&0&0&-4&-4&3&2&-1&-2\\
 0&1&-1&0&0&-2&-2&1&1&0&-1\\
 0&0&0&1&0&-1&0&1&1&0&0\\
 0&0&0&0&1&0&-1&1&1&0&0
 \end{pmatrix}.
\end{equation}
The first nine are the vertices of the reflexive polytope $\Delta_9$;
$v_9$ is interior to a facet and $v_{10}$ is the pentagon centre.
Use the source's integral support vector
\begin{equation}\label{eq:x9-height}
 h=(2,2,2,3,3,3,3,2,2,1,2),\qquad
 P_h=\{u\in\RR^4:h_i+\langle u,v_i\rangle\ge0\text{ for all }i\}.
\end{equation}
Its normal fan is the source fan $\Sigma'$ with twenty-one maximal
cones. The vertices of $P_h$ are
\begin{equation}\label{eq:x9-cartier-vertices}
\begin{array}{rrr}
(-2,-2,-3,15)&(-2,-2,-3,9)&(-2,-2,15,-3)\\
(-2,-2,9,-3)&(-2,-2,15,15)&(-2,4,-3,3)\\
(-2,4,3,-3)&(-2,4,3,3)&(1,-2,-3,3)\\
(0,-2,-3,3)&(1,-2,-3,1)&(1,-2,3,-3)\\
(0,-2,3,-3)&(1,-2,1,-3)&(1,-2,3,3)\\
(0,2,-3,-1)&(0,2,-1,-3)&(0,2,-1,-1)\\
(1,0,-3,-1)&(1,0,-1,-3)&(1,0,-1,-1).
\end{array}
\end{equation}
Substitution in \eqref{eq:x9-height} gives the source cones as the
active sets; every omitted inequality has margin at least one.
The recession cone is zero: pairing with $v_0,v_1,v_2$ first forces
the first two coordinates to vanish, and the remaining opposite
coordinate pairs force the other two to vanish. Thus these integral
vertices give a strictly convex Cartier support function. On
$P'=X_{\Sigma'}$, the divisor $A=\sum h_iD_i$ is ample and Cartier.
Set
\begin{equation}\label{eq:x9-polarization}
 \mathcal L_0=A|_{X'_9},\qquad H=c_1(q^*\mathcal L_0).
\end{equation}
Here $X'_9$ is the proper transform of a sufficiently general
nondegenerate anticanonical member for $\Delta_9$.

For an explicit smooth ambient refinement, replace the heights by
\begin{equation}\label{eq:x9-fine-height}
 a_i=h_i-\frac{\mathbf1_{i=3}+\mathbf1_{i=4}}{100}
             +\frac{i^2}{10^6}.
\end{equation}
The maximal cones of the normal fan of $P_a$ have ray-index sets
\begin{equation}\label{eq:x9-fine-fan}
\begin{gathered}
\{0,1,3,6\},\ \{0,1,3,7\},\ \{0,1,4,5\},\ \{0,1,4,7\},\\
\{0,1,5,6\},\ \{0,2,3,6\},\ \{0,2,3,7\},\ \{0,2,4,5\},\\
\{0,2,4,7\},\ \{0,2,5,6\},\ \{1,3,6,9\},\ \{1,3,7,8\},\\
\{1,3,8,9\},\ \{1,4,5,9\},\ \{1,4,7,8\},\ \{1,4,8,9\},\\
\{1,5,6,9\},\ \{2,3,6,10\},\ \{2,3,7,8\},\ \{2,3,8,10\},\\
\{2,4,5,10\},\ \{2,4,7,8\},\ \{2,4,8,10\},\ \{2,5,6,10\},\\
\{3,6,9,10\},\ \{3,8,9,10\},\ \{4,5,9,10\},\ \{4,8,9,10\},\\
\{5,6,9,10\}.
\end{gathered}
\end{equation}
Each ray matrix has determinant $\pm1$, and each cone lies in a cone
of $\Sigma'$ and in a facet of $\Delta_9$. Consequently its toric
fourfold $T$ is smooth and projective and the refinement is crepant.
At the four source parallelograms
\begin{equation}\label{eq:x9-nodes}
 q_1=(3,6,10,8),\quad q_2=(2,6,3,7),\quad
 q_3=(4,5,10,8),\quad q_4=(2,5,4,7),
\end{equation}
the refinement uses the diagonal $\{a,c\}$ of $q=(a,b,c,d)$.
Thus the anticanonical proper transform $Y\subset T$ is the source's
chosen projective small resolution $q:Y\to X'_9$.
Extra refinements over facet-interior strata do not change the
hypersurface: those strata map to torus-fixed points missed by a
nondegenerate member. In particular the additional nonsimplicial
coarse cone with rays $\{2,7,8,10\}$ is contained in a single original
facet and contributes no fifth node. These model identifications use
\cite[Proposition~3.4]{wuebben-cy-mirror}.

\subsection{A first-order family with the prescribed ample class}
We state explicitly the additional deformation step. The nodal
specializations of CY~V used here formulate the classical criterion;
no mixed-cone extension of the spectral theorem is asserted.

\begin{lemma}\label{lem:polarized-nodal-family}
Let $Z$ be a connected normal projective nodal threefold over $\CC$
with $\omega_Z\simeq\OO_Z$ and
$H^1(\OO_Z)=H^2(\OO_Z)=0$, admitting a projective small crepant
resolution. Suppose the exceptional-curve relation space contains a
vector with no zero coordinate. For every ample Cartier line bundle
$\mathcal L_0$ on $Z$, there is a proper flat projective smoothing
over a disc, with trivial relative canonical bundle and a relatively
ample line bundle restricting to exactly $\mathcal L_0$, such that
every node is smoothed to first order.
\end{lemma}
\begin{proof}
By \cite[Theorem~2.5]{wuebben-cy-smoothing}, the full-support relation
lifts to a global deformation tangent whose image in every local
node-smoothing space is nonzero. The common-volume-form
identifications may rescale equation parameters by nonzero factors,
which preserve this condition. The deformation base is smooth by
\cite[Theorem~4.5]{wuebben-cy-smoothing}; its additional degree-six
hypothesis is vacuous for nodes.

The projective effectivity construction in
Proposition~\ref{prop:eight-family} applies with this chosen tangent.
For completeness, choose $N$ with $\mathcal L_0^N$ very ample and
$H^1(\mathcal L_0^N)=0$. Across a small Artin extension, the
obstruction to lifting this line bundle lies in $H^2(\OO_Z)$ and
the obstruction to lifting its coordinate sections lies in
$H^1(\mathcal L_0^N)$. Both vanish. The lifted map remains a closed
immersion. Embedded deformations therefore map formally smoothly to
abstract deformations, and the Hilbert germ is smooth with a lift of
the selected tangent.

Work on a finite-type Hilbert neighbourhood with integral fibres.
An \'etale neighbourhood supplies a section through a smooth point of
$Z$ \cite[Lemma~37.38.5, Tag~057G]{stacks}. The relative Picard
scheme and its normalized Poincar\'e bundle then exist as in
\cite[Theorems~2.5(2), 4.8 and Exercise~4.3]{kleiman-picard}.
The Picard morphism is smooth at $[\mathcal L_0]$ because
$H^2(\OO_Z)=0$ \cite[Proposition~5.19]{kleiman-picard}.
A further \'etale-local section through that point extends
$\mathcal L_0$ itself. Ampleness persists after shrinking.
Both changes of base induce tangent isomorphisms. A disc with the
chosen tangent in the resulting smooth base retains every nonzero
local derivative. The auxiliary power $N$ neither changes the final
polarization nor reparametrizes the disc.

The equations $\sum z_j^2=\tau_i(t)$ with $\tau_i'(0)\ne0$
make the total space smooth at the nodes. Properness and smoothness
away from them permit shrinking so that the total space and every
noncentral fibre are smooth. The relative dualizing line bundle
restricts to $\omega_Z$. Since $H^1(\omega_Z)=0$, cohomology and
base change lifts a central nowhere-zero generator to a relative
section. Its zero locus has proper image missing the origin; after
shrinking it is empty. This trivializes the relative canonical bundle.
\end{proof}

\begin{proposition}\label{prop:x9-family}
The model $X'_9$, resolution $Y$ and exact polarization
\eqref{eq:x9-polarization} admit a smoothing satisfying all geometric
hypotheses of Theorem~\ref{thm:matrix}. Their second Betti numbers
are $b_2(Y)=6$ and $b_2(X_t)=3$.
\end{proposition}
\begin{proof}
The earlier construction gives exactly four nodes and no other
singularities. Adjunction and crepancy give trivial dualizing sheaves.
On the smooth projective toric fourfold $T$, the anticanonical sequence
$0\to K_T\to\OO_T\to\OO_Y\to0$, toric structure-sheaf
vanishing and Serre duality give $H^1(\OO_Y)=H^2(\OO_Y)=0$.
Rationality of the nodes transfers the vanishings to $X'_9$.
The source's full integral curve-relation lattice is
\begin{equation}\label{eq:x9-curve-relation}
 \left\{\ell\in\ZZ^4:\sum_i\ell_i[C_i]=0\right\}
 =\ZZ(1,-1,-1,1)
\end{equation}
by \cite[Proposition~3.5 and Section~11.1]{wuebben-cy-mirror}.
Lemma~\ref{lem:polarized-nodal-family} therefore supplies the exact
polarized first-order smoothing and trivial relative canonical bundle.

For simple connectivity, apply
\cite[Corollary~1.9]{batyrev-kreuzer} to the nondegenerate hypersurface
and the smooth projective crepant fan above. The relevant lattice
contains all original vertices, in particular the basis
$v_0,v_1,v_3,v_4$ of $\ZZ^4$. Its quotient, hence $\pi_1(Y)$,
is zero. Van Kampen across the simply connected links and fillings
gives $\pi_1(X_t)=0$.

The source divisors span $H^2(Y;\QQ)$
\cite[Lemma~3.3]{wuebben-cy-mirror}. There are four independent
character relations, and $D_9|_Y=0$ because the facet-interior divisor
misses the hypersurface. Thus $b_2(Y)\le6$. The nonsingular
intersection matrix \eqref{eq:x9-intersection} below proves that
six indicated classes are independent. The exceptional span has
dimension three by \eqref{eq:x9-curve-relation}, so the attachment
sequence gives $b_2(X_t)=6-3=3$.
\end{proof}

This constructs a suitable smoothing of the exact earlier central
model. CY~VI's Theorem~9.25 supplies a proper analytic smoothing with
projective fibres; it does not by itself assert the precise relative
polarization or first-order condition used here. We do not identify
the chosen disc with that theorem's arbitrary smoothing arc, or its
metrics with those of a birational mirror model.

\subsection{The three-dimensional harmonic quotient}
Use the rational basis of integral divisor classes
\[
 (D_2,D_5,D_6,D_7,D_8,D_{10}).
\]
No assertion that this is an integral lattice basis is needed.
The character relations, $D_9|_Y=0$ and \eqref{eq:x9-height} give
\begin{equation}\label{eq:x9-h-coordinates}
 H=(6,18,18,-12,-10,8),\qquad H^3=8712.
\end{equation}
The polarized intersection matrix is
\begin{equation}\label{eq:x9-intersection}
 B=\left(\int_Y D_iD_jH\right)=
 \begin{pmatrix}
 -18&6&6&6&0&0\\
 6&0&14&0&0&2\\
 6&14&0&0&0&2\\
 6&0&0&-22&12&0\\
 0&0&0&12&-16&2\\
 0&2&2&0&2&-6
 \end{pmatrix},\qquad \det B=-5064192.
\end{equation}
These are intersection products in the toric ring with character
relations $\sum_i\langle m,v_i\rangle D_i=0$ and the
Stanley--Reisner relations of \eqref{eq:x9-fine-fan}. Every maximal
cone product integrates to one, and $[Y]=\sum_iD_i$.
These relations, together with the displayed rays and cones, determine
every entry of \eqref{eq:x9-intersection} and the volume in
\eqref{eq:x9-h-coordinates}.

For $q=(a,b,c,d)$ with the $ac$ diagonal, the integral source circuit
is $e_b+e_d-e_a-e_c$. In the chosen basis the four circuits give
\begin{equation}\label{eq:x9-period-matrix}
 C=\begin{pmatrix}
 0&0&1&0&1&-1\\
 -1&0&1&1&0&0\\
 0&1&0&0&1&-1\\
 -1&1&0&1&0&0
 \end{pmatrix}.
\end{equation}
It has rank three, with the full surviving kernel spanned by
\begin{equation}\label{eq:x9-kernel}
 k_1=D_2+D_7,\quad k_2=D_5+D_6-D_7+D_{10},\quad
 k_3=D_8+D_{10}.
\end{equation}
In particular $H=6k_1+18k_2-10k_3$. Choose quotient lifts
$(D_2,D_5,D_6)$. The energy and the full Schur complement are
\begin{equation}\label{eq:x9-forms}
 E=\begin{pmatrix}2&-1&-1\\-1&2&0\\-1&0&2\end{pmatrix},
 \qquad
 M=\frac1{15}\begin{pmatrix}
 177&-79&-79\\-79&158&-52\\-79&-52&158
 \end{pmatrix}.
\end{equation}
The leading principal minors of $E$ are $2,3,4$; those of $M$ are
$59/5,869/9,17584/45$. Both forms are positive definite. Direct
expansion gives
\begin{equation}\label{eq:x9-polynomial}
 \frac{\det(E-\mu M)}{-\det M}
 =\left(\mu-\frac17\right)
   \left(\mu^2-\frac{375}{628}\mu+\frac{45}{628}\right).
\end{equation}
The ordered roots are $1/7$ and
$(375\mp3\sqrt{3065})/1256$. They are positive and distinct.
Proposition~\ref{prop:x9-family} verifies the geometric assumptions,
and Theorem~\ref{thm:matrix} now proves
Corollary~\ref{cor:x9}, including its multiplicities, type and gaps.

The earlier vector $(1,-1,-1,1)$ is the kernel of the curve-class
map $A$ and provides a full-support smoothing tangent. Its annihilator
is the three-dimensional period image of $C=A^*$ and counts the
small modes. The new coefficients additionally use the global form
\eqref{eq:x9-intersection} and projection off all three classes in
\eqref{eq:x9-kernel}. The mirror equality of relation lattices in
CY~VI is not used to calculate that norm. Thus the earlier construction
and deformation theorems are actual premises of this application,
while the quantitative spectral conclusion follows from the analysis
of the present paper.
